%============================
% CHAPTER 5: WEBGL RENDERING
%============================
\section{WebGL Rendering: GPU-Powered Visualization}

\begin{frame}
\frametitle{Chapter 5: Awakening the GPU}

\begin{columns}[T]
\begin{column}{0.6\textwidth}
  \textbf{Dr. Chen's Realization:} "JavaScript alone can't render 50 million voxels at 60 FPS. I need the GPU!"
  
  \vspace{0.5cm}
  \textbf{WebGL} is the bridge between JavaScript and GPU hardware:
  \begin{itemize}
    \item \textbf{JavaScript}: Handles data loading and logic
    \item \textbf{WebGL API}: GPU resource management
    \item \textbf{Shaders}: GPU programs for rendering
    \item \textbf{Canvas}: Display output
  \end{itemize}
\end{column}
\begin{column}{0.35\textwidth}
  \begin{center}
    \begin{tikzpicture}[scale=0.8]
      % GPU pipeline
      \draw[fill=medicalblue!20, rounded corners] (0,0) rectangle (2,0.5);
      \node[medicalblue] at (1,0.25) {JavaScript};
      
      \draw[fill=softgreen!30, rounded corners] (0,0.7) rectangle (2,1.2);
      \node[medicalgreen] at (1,0.95) {WebGL API};
      
      \draw[fill=medicalred!20, rounded corners] (0,1.4) rectangle (2,1.9);
      \node[medicalred] at (1,1.65) {GPU Shaders};
      
      \draw[fill=softblue!30, rounded corners] (0,2.1) rectangle (2,2.6);
      \node[medicalblue] at (1,2.35) {Canvas Display};
      
      % Arrows
      \draw[->, thick] (1,0.5) -- (1,0.7);
      \draw[->, thick] (1,1.2) -- (1,1.4);
      \draw[->, thick] (1,1.9) -- (1,2.1);
    \end{tikzpicture}
    \vspace{0.3cm}
    {\small\color{medicalgray}WebGL Pipeline}
  \end{center}
\end{column}
\end{columns}

\knowledgecheck{GPU vs CPU}
Why can't we render medical volumes using only JavaScript? What specific advantages does the GPU provide?
\end{frame}

\begin{frame}{WebGL Pipeline: GPU Resources}
\begin{center}
\begin{tikzpicture}[
    node distance=1cm,
    res/.style={rectangle, rounded corners, draw=medicalblue!60!black, fill=softblue!20, thick, text width=2.8cm, align=center, font=\sffamily\small, minimum height=1cm}
]
  \node[res] (tex) {3D Texture \\ \scriptsize (Voxel Data)};
  \node[res, left=of tex] (buf) {Geometry Buffers \\ \scriptsize (Quad Vertices)};
  \node[res, right=of tex] (shd) {Shader Source \\ \scriptsize (GLSL Code)};

  \draw[dashed, medicalgray] (-5, -1) -- (5, -1) node[midway, below, font=\tiny] {Uploaded from CPU to GPU VRAM};
\end{tikzpicture}
\end{center}

\begin{block}{Volume Storage}
The entire 3D DICOM stack is stored as a \textbf{Texture3D}, allowing the GPU to perform hardware-accelerated interpolation between slices.
\end{block}

\clinicalexample{GPU Memory}
A 512×512×200 CT volume requires ~200MB of GPU memory. Modern GPUs can handle this, but memory management is critical.
\end{frame}

\begin{frame}{WebGL Pipeline: Shader Execution}
\begin{center}
\begin{tikzpicture}[
    node distance=1.2cm,
    pipe/.style={rectangle, rounded corners, draw=medicalgreen!60!black, fill=softgreen!20, thick, text width=3.5cm, align=center, font=\sffamily\small, minimum height=0.9cm},
]
  \node[pipe] (vs) {Vertex Shader \\ \scriptsize (Positions the Quad)};
  \node[pipe, below=of vs] (fs) {Fragment Shader \\ \scriptsize (Samples Texture)};
  \node[pipe, below=of fs] (canvas) {Canvas Display \\ \scriptsize (Final Output)};

  \draw[-Stealth, thick, medicalgray] (vs) -- (fs);
  \draw[-Stealth, thick, medicalgray] (fs) -- (canvas);
\end{tikzpicture}
\end{center}

\begin{itemize}
    \item \textbf{Rasterization:} The GPU determines which pixels on screen cover the geometry (our slice quad).
    \item \textbf{Texture Sampling:} For every fragment, the shader samples the 2D texture to get the grayscale value.
\end{itemize}

\knowledgecheck{Shader Roles}
What does the vertex shader do? What does the fragment shader do? Why do we need both?
\end{frame}

\begin{frame}[fragile]
\frametitle{WebGL Context Initialization}

\begin{lstlisting}
constructor(canvas: HTMLCanvasElement, 
            vertexSource: string, 
            fragmentSource: string) {
  this.canvas = canvas;
  
  // Get WebGL context
  const gl = canvas.getContext('webgl');
  if (!gl) throw new Error('WebGL not supported');
  
  this.gl = gl;
  
  // Enable depth testing
  this.gl.enable(this.gl.DEPTH_TEST);
  this.gl.depthFunc(this.gl.LEQUAL);
  
  // Set viewport
  this.gl.viewport(0, 0, canvas.width, canvas.height);
  
  // Compile shaders
  this.program = new ShaderProgram(gl, 
                                   vertexSource, 
                                   fragmentSource);
  
  // Create camera and controls
  this.camera = new Camera();
  this.mouse = new MouseController(canvas, this.camera);
}
\end{lstlisting}
\end{frame}

\begin{frame}
\frametitle{Shader Program Compilation}

\textbf{Pipeline:}
\begin{enumerate}
  \item \textbf{Create Shader}: \texttt{gl.createShader(type)}
  \item \textbf{Set Source}: \texttt{gl.shaderSource(shader, source)}
  \item \textbf{Compile}: \texttt{gl.compileShader(shader)}
  \item \textbf{Create Program}: \texttt{gl.createProgram()}
  \item \textbf{Attach Shaders}: \texttt{gl.attachShader(program, shader)}
  \item \textbf{Link}: \texttt{gl.linkProgram(program)}
  \item \textbf{Use}: \texttt{gl.useProgram(program)}
\end{enumerate}

\vspace{0.5cm}

\textbf{Error Handling:}
\begin{itemize}
  \item Check \texttt{gl.COMPILE\_STATUS} after compilation
  \item Check \texttt{gl.LINK\_STATUS} after linking
  \item Use \texttt{getShaderInfoLog()} for debug messages
\end{itemize}
\end{frame}

\begin{frame}[fragile]
\frametitle{Vertex Shader Example}

\textbf{Purpose:} Transform geometry to screen space

\begin{lstlisting}[language=glsl]
attribute vec3 a_position; // Vertex position
attribute vec2 a_texcoord; // Texture coordinate
varying vec2 v_texcoord;   // To fragment shader
uniform mat4 u_matrix;     // Model-view-projection

void main() {
  // Pass texture coordinate to fragment shader
  v_texcoord = a_texcoord;
  // Transform position
  gl_Position = u_matrix * vec4(a_position, 1.0);
}
\end{lstlisting}

\textbf{Key Concepts:}
\begin{itemize}
  \item \textbf{attribute}: Per-vertex data
  \item \textbf{uniform}: Constant for all vertices
  \item \textbf{varying}: Interpolated per fragment
  \item \textbf{gl\_Position}: Output vertex position
\end{itemize}
\end{frame}

\begin{frame}[fragile]
\frametitle{Fragment Shader Example}

\textbf{Purpose:} Determine color for each pixel

\begin{lstlisting}[language=glsl]
precision mediump float;
varying vec2 v_texcoord;
uniform sampler2D u_texture;

void main() {
  // Sample texture at this fragment
  float v = texture2D(u_texture, v_texcoord).r;
  
  // Convert to grayscale color
  gl_FragColor = vec4(v, v, v, 1.0);
}
\end{lstlisting}

\textbf{Key Concepts:}
\begin{itemize}
  \item \textbf{sampler2D}: 2D texture unit
  \item \textbf{texture2D()}: Sample texture at UV coords
  \item \textbf{gl\_FragColor}: Output pixel color
  \item \textbf{precision}: Floating-point precision
\end{itemize}
\end{frame}

\begin{frame}[fragile,allowframebreaks]
\frametitle{Textures in WebGL}

\textbf{Texture Creation:}
\begin{itemize}
  \item \texttt{gl.createTexture()}: Allocate texture object
  \item \texttt{gl.bindTexture(gl.TEXTURE\_2D, tex)}: Bind for operations
  \item \texttt{gl.texImage2D()}: Upload pixel data
  \item \texttt{gl.texParameteri()}: Set filtering/wrapping
\end{itemize}

\vspace{0.3cm}

\textbf{Slice Texture Update:}
\begin{lstlisting}[language=JavaScript, basicstyle=\small\ttfamily]
// Create texture from slice data
gl.bindTexture(gl.TEXTURE_2D, this.sliceTexture);
gl.texImage2D(
  gl.TEXTURE_2D, 0, gl.LUMINANCE,
  width, height, 0,
  gl.LUMINANCE, gl.UNSIGNED_BYTE,
  sliceBytes
);
\end{lstlisting}

\vspace{0.3cm}

\textbf{Key Parameters:}
\begin{itemize}
  \item \textbf{LUMINANCE}: Single channel (grayscale). Perfect for DICOM!
  \item \textbf{UNSIGNED\_BYTE}: 8-bit values ($0-255$).
  \item \textbf{FILTERING}: \texttt{gl.LINEAR} allows the GPU to interpolate between pixels for smooth zooming.
\end{itemize}
\end{frame}

\begin{frame}
\frametitle{Vertex Buffers and Attributes}

\textbf{Vertex Buffer Object (VBO):}
\begin{itemize}
  \item GPU memory object storing vertex data
  \item \texttt{gl.createBuffer()}: Create buffer
  \item \texttt{gl.bindBuffer()}: Bind for operations
  \item \texttt{gl.bufferData()}: Upload data to GPU
\end{itemize}

\vspace{0.5cm}

\textbf{Vertex Attribute Pointers:}
\begin{itemize}
  \item \texttt{gl.enableVertexAttribArray()}: Enable attribute
  \item \texttt{gl.vertexAttribPointer()}: Describe layout
  \item \textbf{Parameters}: Index, size, type, normalized, stride, offset
\end{itemize}

\vspace{0.5cm}

\textbf{Example:} Position (vec3) + TexCoord (vec2)
\begin{itemize}
  \item Stride: $3 \times 4 + 2 \times 4 = 20$ bytes
  \item Offset for texcoord: $3 \times 4 = 12$ bytes
\end{itemize}
\end{frame}

\begin{frame}[fragile,allowframebreaks]
\frametitle{Rendering Loop}

\begin{lstlisting}
start(): void {
  const loop = () => {
    // Clear buffers
    this.gl.clear(this.gl.COLOR_BUFFER_BIT | 
                  this.gl.DEPTH_BUFFER_BIT);
    
    // Render scene
    if (this.mesh) {
      this.program.use();
      this.renderMesh(this.mesh);
    }
    
    // Render slices
    if (this.volume) {
      this.sliceProgram.use();
      this.renderSlices();
    }
    
    // Request next frame
    this.animationId = requestAnimationFrame(loop);
  };
  
  this.animationId = requestAnimationFrame(loop);
}
\end{lstlisting}

\textbf{Key Points:}
\begin{itemize}
  \item \textbf{requestAnimationFrame}: Browser-optimized rendering
  \item \textbf{Clear}: Remove previous frame data
  \item \textbf{Use}: Activate shader program
  \item \textbf{Draw}: Render geometry
\end{itemize}
\end{frame}

\begin{frame}
\frametitle{Depth Testing}

\textbf{Problem:} Overlapping geometry renders incorrectly without depth info

\vspace{0.5cm}

\textbf{Solution:} Depth Test (Z-buffering)
\begin{itemize}
  \item Each pixel stores depth value (distance from camera)
  \item Only draw pixel if new depth $<$ stored depth
  \item Ensures correct occlusion ordering
  \item \texttt{gl.enable(gl.DEPTH\_TEST)}
  \item \texttt{gl.depthFunc(gl.LEQUAL)}: Update if less or equal
\end{itemize}

\vspace{0.5cm}

\textbf{Configuration:}
\begin{itemize}
  \item \texttt{NEVER}: Never pass
  \item \texttt{LESS}: Pass if new $<$ stored
  \item \texttt{EQUAL}: Pass if equal
  \item \texttt{LEQUAL}: Pass if $\leq$ (typical)
  \item \texttt{GREATER}: Pass if $>$
  \item \texttt{ALWAYS}: Always pass
\end{itemize}
\end{frame}

\begin{frame}[fragile]
\frametitle{Matrix Transformations}

\textbf{Model-View-Projection (MVP) Matrix:}
$$\text{gl\_Position} = P \times V \times M \times \text{position}$$

Where:
\begin{itemize}
  \item $M$: Model matrix (geometry to world space)
  \item $V$: View matrix (world to camera space)
  \item $P$: Projection matrix (camera to screen space)
\end{itemize}

\vspace{0.5cm}

\textbf{Camera Implementation:}
\begin{itemize}
  \item Spherical coordinates: azimuth, polar, distance
  \item Eye position: $(r \sin\phi \cos\theta, r \sin\phi \sin\theta, r \cos\phi)$
  \item Look-at: Compute view matrix from eye, center, up
  \item Perspective: Configurable FOV, aspect ratio, near/far planes
\end{itemize}
\end{frame}